\section{Gaussian hierarchies and their calculus}\label{sec:foundations}

We construct the probability spaces and calculus for the next two sections.
Pressure differentiation turns covariance perturbations into overlap
errors in Section~\ref{sec:comparisons}; moments on one planted path
control those errors, and the finite-tree calculus organizes the
fixed-degree cavity expansions of Section~\ref{sec:cavity}.
Throughout this section
$n$ is a positive integer and the prior on $\{-1,1\}^n$ is uniform. All
expectations include both the Gaussian variables and the indicated spin
draws. The normalized planted law and its deterministic-weight variance
estimate follow the framework of \cite[Section~2 and Proposition~2.2]{AL2026}.
We prove the higher-moment and finite-tree extensions needed here.

\subsection{Finite hierarchies and conditional covariances}

Put $y(\sigma)=\sigma/\sqrt n$ and
$R(\sigma,\tau)=y(\sigma)\cdot y(\tau)$. The two feature maps
\[
 \phi_b(\sigma)=\frac{\sigma\sigma^{\mathsf T}}{\sqrt{2n}},
 \qquad \phi_h(\sigma)=\sigma
\]
take values in Euclidean spaces of dimensions $n^2$ and $n$, respectively.
Their inner products are $nR^2/2$ and $nR$, and their squared norms are
$n/2$ and $n$. The use of all ordered spin products is harmless and makes
the covariance realization immediate.

For the moment let $b,h$ be nonnegative nondecreasing step functions on
$(0,1)$. Choose a common partition
\[
 0=m_0<m_1<\cdots<m_L<m_{L+1}=1,
 \qquad b(u)=b_j,\quad h(u)=h_j
 \quad(m_j\leq u<m_{j+1}).
\]
Set $b_{-1}=h_{-1}=0$. Write $Z_j$ for independent feature-space Gaussian
vectors having coordinate variances $b_j-b_{j-1}$ in the first feature
space and $h_j-h_{j-1}$ in the second. Zero variances are permitted.
For a combined external field $z=(z_b,z_h)$ define
\[
 F_L(z)=\log\left(2^{-n}\sum_\sigma
       e^{z_b\cdot\phi_b(\sigma)+z_h\cdot\phi_h(\sigma)}\right),
 \qquad
 F_{j-1}(z)=\frac1{m_j}\log\mathbf E e^{m_jF_j(z+Z_j)}
 \quad (j=L,\ldots,1),
\]
and $F_{-1}(0)=\mathbf E F_0(Z_0)$. Thus the operation at mass zero is
expectation. Set
\[
 f_n(b,h)=\frac1n F_{-1}(0)-\frac{b_L}{4}-\frac{h_L}{2}.
\]
Common linear spin sources may be included in every terminal partition
function. A source coupled to $y$, used only for differentiation and then
set to zero, will be called a normalized spin source.

The associated planted law is constructed forwards. Draw $Z_0$ from its
Gaussian law. Conditional on $G_{j-1}=\sum_{i<j}Z_i$, draw $Z_j$ with
density
\[
 \exp\{m_j[F_j(G_{j-1}+Z_j)-F_{j-1}(G_{j-1})]\}
\]
relative to its Gaussian law. Finally draw $\sigma$ from the Gibbs law
in the field $G_L$. These are normalized transition kernels. At a
specified split, replicas share the prefix and then use conditionally
independent copies of all subsequent transitions. A finite deterministic
tree of such splits defines a probability law, even when some split
masses coincide or some covariance increments vanish. At a jump we
specify whether the split is before or after the increment. Repeated
empty steps allow these conventions to be recorded in a common partition.

Let $\mathcal F_u$ be the prefix sigma-field after the revelation specified
at $u$. For $m_j\leq u<m_{j+1}$ it includes $G_j$. Define
\[
 \begin{gathered}
 X_u=\mathbf E[y\mid\mathcal F_u],\qquad
 C_u=\mathbf E[yy^{\mathsf T}\mid\mathcal F_u]-X_uX_u^{\mathsf T},\\
 S_u=\mathbf E|X_u|^2,\qquad
 B_u=\mathbf E\bigl\|\mathbf E[yy^{\mathsf T}\mid\mathcal F_u]\bigr\|_{\rm HS}^2,
 \qquad D_u^2=B_u-S_u^2,\\
 L_u=y\cdot X_u-\tfrac12|X_u|^2.
 \end{gathered}
\]
Here and below conditional expectation given a prefix means conditioning
on its entire sigma-field.

\begin{lemma}[Source derivatives]\label{lem:found-source}
The gradient of a continuation in a source coupled to a bounded feature
$\phi$ is the conditional mean of $\phi$. Its Hessian is the sum of the
remaining conditional martingale covariance increments, weighted by the
mass of each transition, with mass $1$ for the final spin draw. In
particular, the Hessian $H_u$ in a normalized spin source satisfies
\[
 uC_u\preceq H_u\preceq C_u.
\]
If $y_1,y_2$ are independent continuations from the prefix at $u$, then
\[
 \mathbf E(y_1\cdot y_2)=S_u,\qquad
 \mathbf E(y_1\cdot y_2)^2=B_u,
\]
and hence $D_u$ is the standard deviation of this overlap. Moreover,
\[
 D_u^2=\operatorname{Var}(|X_u|^2)
       +2\mathbf E[X_u^{\mathsf T}C_uX_u]
       +\mathbf E\operatorname{tr}C_u^2.
\]
\end{lemma}

\begin{proof}
For $\mathcal T_m F=m^{-1}\log\mathbf E e^{mF}$, differentiation in any
external source gives
\[
 \nabla\mathcal T_mF=\mathbf E_m\nabla F,\qquad
 \nabla^2\mathcal T_mF=\mathbf E_m\nabla^2F
                  +m\operatorname{Cov}_m(\nabla F),
\]
where $\mathbf E_m$ is the normalized tilted expectation. At the
terminal spin sum the same formulas hold with mass $1$ and gradient
$\phi$. Iteration proves the claimed Hessian representation. More
explicitly, if $M_j=\mathbf E[\phi\mid\mathcal F_{m_j}]$ and
$C_L^\phi=\operatorname{Cov}(\phi\mid\mathcal F_{m_L})$, the Hessian at
the prefix indexed by $j$ is
\[
 \mathbf E\left[C_L^\phi+
   \sum_{k=j+1}^L m_k(M_k-M_{k-1})(M_k-M_{k-1})^{\mathsf T}
           \,\middle|\,\mathcal F_{m_j}\right].
\]
The same expression without the mass factors is the conditional
covariance of $\phi$. All remaining masses belong to $[u,1]$.
This proves the matrix inequalities. Conditional independence proves the
two overlap identities. Finally expand
$\|C_u+X_uX_u^{\mathsf T}\|_{\rm HS}^2$ and subtract $S_u^2$.
\end{proof}

\subsection{Pressure derivatives and the pinned moment estimate}

The three estimates needed later can be collected as follows. In the
first line, $db(u),dh(u)$ are variations in a parameter indexing paths;
they are not Stieltjes increments in $u$. In the second line $r$ increases
along the planted path. Its derivative is taken on absolutely continuous
parts. At a jump of the pair of clocks, fix the joint filling
$(b_-+t\Delta b,h_-+t\Delta h)$, $0\le t\le1$, at the jump's constant
mass. All intermediate revelations refer to this chosen filling.
The pressure and susceptibility estimates allow a common deterministic
source. The moment estimate is at zero deterministic external field,
under the full unconditioned law of one planted path and its terminal
spin configuration.
For every deterministic $k\in L^2(0,1)$ and $2\leq p<\infty$,
\begin{equation}
\begin{split}
 df_n&=-\tfrac14\int_0^1B_u\,db(u)\,du
       -\tfrac12\int_0^1S_u\,dh(u)\,du,\\
 \partial_rS_u&\geq n\,\partial_rh\,
                  \mathbf E\operatorname{tr}H_u^2,\\
 \left\|\int_0^1(L_u-\mathbf EL_u)k_u\,du\right\|_p^2
     &\leq c_p\left\|\int_0^1 k_u^2W_u\,du\right\|_{p/2},
 \qquad W_u=(y-X_u)^{\mathsf T}H_u(y-X_u),\\[-2pt]
 &\hspace{100pt}\mathbf EW_u=\mathbf E\operatorname{tr}(H_uC_u).
\end{split}\tag{1}\label{eq:1}
\end{equation}
One may take $c_2=1$ and $c_p\leq1+p^2/3$.

\begin{lemma}[Gaussian pressure differentiation]\label{lem:found-pressure}
The first two lines of \eqref{eq:1} hold for finite hierarchies, including
degenerate Gaussian increments by continuity. In addition, for any common
external source and any two admissible step paths,
\[
 |f_n(b,h)-f_n(\widetilde b,\widetilde h)|
 \leq\tfrac14\|b-\widetilde b\|_{L^1}
        +\tfrac12\|h-\widetilde h\|_{L^1}.
\]
The bound is uniform in the source when the same self correction is used.
\end{lemma}

\begin{proof}
Consider first one feature map $\phi$ of constant squared norm $q$ and
its cumulative covariance clock $c_j$. Let
$A_j=\mathbf E|\mathbf E[\phi\mid\mathcal F_{m_j}]|^2$.
Gaussian heat differentiation of the increment $c_k-c_{k-1}$, followed by
the tilted chain rule through all earlier transitions, gives
\[
 \frac{\partial F_{-1}(0)}{\partial(c_k-c_{k-1})}
    =\tfrac12\mathbf E\{\operatorname{tr}H_k^\phi
                         +m_k|M_k|^2\}.
\]
For $k=0$ the second term is zero. The martingale decomposition in
Lemma~\ref{lem:found-source} gives the explicit telescoping sum
\[
 \mathbf E\operatorname{tr}H_k^\phi
     =q-A_L+\sum_{j=k+1}^L m_j(A_j-A_{j-1}),
\]
so the last derivative equals
\[
 \tfrac q2-\tfrac12\sum_{j=k}^L(m_{j+1}-m_j)A_j.
\]
Subtract $qc_L/2$ from $F_{-1}(0)$. On writing
$d(c_k-c_{k-1})=dc_k-dc_{k-1}$, the derivative becomes
$-\frac12\sum_j(m_{j+1}-m_j)A_j\,dc_j$.
For the matrix feature $A_j=nB_{m_j}/2$; for the field feature
$A_j=nS_{m_j}$. Dividing by $n$ proves the pressure formula. All spin
bounds used here remain true with an external source. Since $0\leq S,B\leq1$,
integration along the convex segment joining two monotone paths proves
the asserted Lipschitz bound.

To prove the second line, fill one Gaussian increment continuously at its
fixed mass $m$. Its continuation solves the backwards heat equation with
nonlinearity $m|\nabla F|^2/2$. Under the normalized planted transition the
driving Gaussian process acquires drift $m\nabla F$. Differentiating the
backwards equation cancels this drift in the stochastic differential of
the source gradient. Thus $X$ is a martingale throughout the filled
increment. Its diffusion coefficient against the spin field of clock
$h$ is $\sqrt n H$, since the actual field couples to $\sqrt n y$.
The independent matrix field contributes another positive semidefinite
quadratic variation. It\^o's formula therefore gives
\[
 d\mathbf E|X|^2
    =n\mathbf E\operatorname{tr}H^2\,dh
       +\mathbf E\|J_b\|_{\rm HS}^2\,db,
\]
where $J_b$ is the mixed Hessian for the matrix feature and normalized
spin source. The last term is nonnegative. Concatenation proves the
claim along the path and through filled jumps. The heat differentiations
may first be made with positive variances and then passed to zero
variance by dominated convergence of the finite Gaussian integrals.
\end{proof}

\begin{lemma}[Pinned moments for one planted path]\label{lem:found-moments}
At zero normalized spin source and with no other deterministic external
field, the last two lines of \eqref{eq:1} hold with $c_2=1$ and
$c_p\leq1+p^2/3$. All these moments concern one planted path and its
terminal spin, with the full, unconditioned prefix law.
\end{lemma}

\begin{proof}
Write all independent Gaussian increments as linear images of a standard
Gaussian vector $\xi$. Multiply their normalized transition densities
and the terminal Gibbs density, and then condition on the entire terminal
configuration $\sigma\in\{-1,1\}^n$.
Terms at neighboring levels cancel. Up to an additive constant, the
negative logarithm of the resulting density of $\xi$ is
\[
 U_\sigma(\xi)=\tfrac12|\xi|^2-\ell_\sigma(\xi)
                      +\sum_{j=0}^L(m_{j+1}-m_j)F_j(G_j),
\]
where $\ell_\sigma$ is linear. Each continuation is convex, by
Lemma~\ref{lem:found-source}. If $A_u$ is its Hessian in the underlying
drivers, then
\[
 Q:=\nabla^2U_\sigma=I+\int_0^1 A_u\,du\succeq I.
\]
Use a virtual normalized spin source $\lambda$. The joint Hessian of
the continuation in $(\xi,\lambda)$ has blocks
\[
 \begin{pmatrix}A_u&B_u^{\rm mix}\\
                 (B_u^{\rm mix})^{\mathsf T}&H_u\end{pmatrix}
 \succeq0.
\]
With the spin fixed, the driver gradient of $L_u$ is
$g_u=B_u^{\rm mix}(y-X_u)$. Positivity of this block matrix implies, for
every driver direction $v$,
\[
 |v\cdot g_u|^2\leq(v^{\mathsf T}A_uv)W_u.
\]
Put $G=\int_0^1k_ug_u\,du$ and $J=\int_0^1k_u^2W_u\,du$.
Cauchy--Schwarz in $u$ gives
\[
 |v\cdot G|^2\leq
   \left(\int_0^1v^{\mathsf T}A_uv\,du\right)J
   \leq(v^{\mathsf T}Qv)J.
\]
Taking $v=Q^{-1}G$ proves $G^{\mathsf T}Q^{-1}G\leq J$.

The classical Brascamp--Lieb variance inequality
\cite[Theorem~4.1]{BL1976} applies to
this finite-dimensional density: $U_\sigma$ is smooth, its Hessian is
bounded below by $I$, and it has Gaussian tails after a linear tilt.
All required derivatives have at most the integrable growth allowed by
this density; alternatively, one can first use smooth compactly supported
approximations. It yields the conditional variance bound by
$\mathbf E[J\mid\sigma]$.

There is no change of centering on removing the pin. For
$\eta\in\{-1,1\}^n$, the transformation
\[
 \sigma_i\mapsto\eta_i\sigma_i,\quad
 (Z_h)_i\mapsto\eta_i(Z_h)_i,\quad
 (Z_b)_{ij}\mapsto\eta_i\eta_j(Z_b)_{ij}
\]
is orthogonal in each Gaussian feature space, preserves the hierarchy,
and sends $X$ to $\eta X$. Consequently it preserves $L_u$ and $W_u$.
This group acts transitively on terminal spins. The conditional law of
$\int k_uL_u\,du$, as well as that of $J$, is therefore independent of the
pinned spin. In particular its conditional mean is its unconditional
mean. Brascamp--Lieb proves the assertion for $p=2$ with $c_2=1$.

Here is an explicit extension to all real $p>2$. Set
$A=\int k_u(L_u-\mathbf EL_u)\,du$. Apply the same variance inequality
to $|A|^{p/2}$, using smooth regularizations at zero when necessary.
H\"older's inequality gives
\[
 \|A\|_p^p\leq\|A\|_{p/2}^p
       +\frac{p^2}{4}\|A\|_p^{p-2}\|J\|_{p/2}.
\]
If $2<p\leq4$, the first term is bounded using
$\|A\|_{p/2}\leq\|A\|_2$. For $p>4$ use the previously established
bound at exponent $p/2$ and
$\|J\|_{p/4}\leq\|J\|_{p/2}$. If $\|J\|_{p/2}=0$, the variance bound
already gives $A=0$. Otherwise put
$t=\|A\|_p^2/\|J\|_{p/2}$ and $\alpha=p/2$. Then
\[
 t^\alpha\leq d^\alpha+(p^2/4)t^{\alpha-1},
 \qquad d=1\ (p\leq4),\quad d=c_{p/2}\ (p>4).
\]
This implies $t\leq d+p^2/4$: for $t\geq d$ divide by
$t^{\alpha-1}$, and for $t<d$ the claim is immediate. Induction down the
finite sequence $p,p/2,p/4,\ldots$ now gives $c_p\leq1+p^2/3$.
Finally, conditioning on the prefix gives
$\mathbf EW_u=\mathbf E\operatorname{tr}(H_uC_u)$.
All integrals are justified first for bounded $k$ and then by $L^2$
approximation; $|X|\leq1$ and $0\preceq H\preceq C\preceq I$ suffice
for this passage.
\end{proof}

\begin{remark}\label{rem:found-single-pin}
For each fixed enlarged tree, first marginalize the other branches and
apply Lemma~\ref{lem:found-moments} to the resulting single-path law.
Then combine its moment bounds by unconditional H\"older under the
positive tree law, and finally integrate the signed topology measure.
The lemma is not applied after spin-dependent reweighting or conditioning
on a fixed prefix; no joint-spin pinned inequality is used.
\end{remark}

\subsection{Signed replica sums at real masses}

The following convention permits covariance differentiation without any
integer-replica or analytic-continuation assumption. An observed tree has
a fixed finite set of labeled leaves, deterministic split masses, and
specified revelation sides. A new summation label is allowed to coincide
with an old leaf, to follow it for a time and then branch, or to branch
at an existing fork. The corresponding expectations are always ordinary
normalized planted expectations. Only the coefficients with which they
are summed may be signed.

\Needspace{3\baselineskip}
\begin{proposition}[Covariance differentiation]\label{prop:found-replica}
Let $K_\theta(u;\sigma,\tau)$ be an affine family of cumulative Gaussian
covariance kernels of a finite hierarchy. Assume that the independent
increment covariances are positive semidefinite throughout the parameter
interval. Let $V=\partial_\theta K_\theta$ and let $P$ be a fixed bounded
spin test on a fixed observed split tree. Then
\begin{equation}
 \frac{d}{d\theta}\nu_\theta[P]
      =\frac12\sum_{c,d}\nu_\theta[P V_{cd}].
 \tag{2}\label{eq:2}
\end{equation}
The signed sum includes coincidences with old labels and between $c,d$.
For distinct labels, $V_{cd}$ is evaluated at their common revelation;
on the diagonal it is evaluated at the terminal covariance. The exact
weights are the real-power weights described in the proof. A constant
diagonal contribution, independent of the summation label, sums to zero.
Conditional means in $P$ can be represented by independent additional
replicas at their specified prefixes before this formula is applied.
\end{proposition}

\begin{proof}
It is enough to prove the formula for a finite hierarchy with positive
masses and nondegenerate increments; degenerate cases follow by
one-sided limits. At a node with incoming mass $a$ and next child mass
$b$, let $I$ denote the next child integral. Without observations the
expression at this node is $I^{a/b}$. If $k$ distinguished children are
prescribed, normalization of the planted transitions gives precisely
\[
 I^{a/b-k}\prod_{i=1}^k I_i,
\]
where $I_i$ is the corresponding tested child integral. A general spin
test is a finite linear combination of products of single-slot tests, so
this identity also determines its expectation by linearity. This finite
expansion is used only for the identity, not to estimate its norm.

Gaussian differentiation at one covariance increment supplies two
derivative slots and a factor $1/2$. Above that increment ordinary chain
rules propagate the integral derivative; below it product and power
rules allocate the two field derivatives. A hit on an existing child
has multiplicity one. If $h$ distinct new children are used, the
power-rule coefficient is $(a/b-k)^{\underline h}$. Each new child
integral carries its incoming factor $b$, so its effective coefficient
is the polynomial
\[
 b^h(a/b-k)^{\underline h}
       =\prod_{j=0}^{h-1}\bigl(a-(k+j)b\bigr).
\]
This identity holds for the actual real masses. In particular it has no
division by a small mass. A double hit in the same child is retained as
such. The final spin summation is the same product calculation with
child mass $1$.

More explicitly, give every derivative slot a distinct label in a finite
set $S$, even when directions agree, and write $D_A=\prod_{s\in A}D_s$,
$D_\varnothing=1$, and $\alpha=a/b-k$. The labeled product rule is
\[
 D_S\!\left(I^\alpha\prod_{i=1}^k I_i\right)
 =\sum_{(S_1,\ldots,S_k,\pi)}
   \alpha^{\underline{|\pi|}}I^{\alpha-|\pi|}
   \prod_{B\in\pi}D_BI\prod_{i=1}^kD_{S_i}I_i.
\]
Here the $S_i$ are disjoint subsets assigned to existing children, and
$\pi$ is an unordered partition of their complement in $S$ into nonempty
blocks, one per new child. Thus no additional $|\pi|!$ occurs. Factoring
the incoming $b$ from every new-child derivative gives the polynomial
coefficients above; zero masses use these polynomial coefficients and
the one-sided limits of the normalized child laws. Recursing with these
labeled partitions proves equality of each final labeled-topology
coefficient for every order of allocating the slots, before any spin
expectation is evaluated. Coincident labels occupy one terminal node.
Consequently allocation order may be changed with any bounded function
of the final topology carried along, before taking absolute variation.

For completeness, at the root perform the calculation on
$e^{a_0F_{\rm root}}\nu[P]$ with an auxiliary incoming power $a_0$.
Set $a_0=0$ after differentiation. The derivative of the prefactor then
vanishes and the normalized expectation remains. This is an algebraic
identity in derivatives of positive real powers; there is no inference
from values at integer masses. Summing the changes of independent
covariance increments assigns to a pair the change accumulated along
exactly its shared prefix, which is $V_{cd}$. This proves \eqref{eq:2}.

The one-label sum at the root has total weight $a_0=0$, as can also be
seen directly in Lemma~\ref{lem:found-attachment}. Hence a constant
diagonal term multiplied by an old-label test sums to zero. A common
source differentiation uses one derivative slot in the identical
calculation. Auxiliary leaves representing conditional means may be
marginalized at any point because all their transitions are normalized.
\end{proof}

\begin{lemma}[Attachment measures and row sums]\label{lem:found-attachment}
Fix an enclosing block of mass $a\in[0,1]$ containing $r\geq1$ observed
leaves. After grouping placements which induce the same extended tree,
one free summation label has the following signed attachment measure:
\begin{enumerate}[label=(\roman*)]
\item weight $+1$ for coincidence with each observed leaf;
\item measure $-dm$ on every occupied open edge, parameterized by its
      increasing mass coordinate;
\item weight $(1-d_v)u_v$ at an observed fork of mass $u_v$ with $d_v$
      occupied children, including an initial fork if present.
\end{enumerate}
The total signed mass is $a$ and the total variation is $2r-a\leq2r$.
For an empty block the first label has weight $a$. Consequently,
adjoining $k$ labels to a tree with $r\geq1$ old labels has total
variation at most
\[
 2^k r(r+1)\cdots(r+k-1).
\]
Adjoining $k\geq1$ labels in a previously empty block has variation at
most $a\,2^{k-1}(k-1)!$. These bounds are independent of the mass
partition. A free row whose summand is independent of the detailed
attachment inside a block of mass $u$ collapses exactly to $u$ times that
summand.
\end{lemma}

\begin{proof}
At a node of mass $a'$ having $d$ occupied children of next mass $b'$,
the coefficient of a previously unoccupied child is $a'-db'$, by the
one-slot real-power rule. If $d=1$, consecutive refinements telescope to
the negative of the length of the occupied edge. When several occupied
children first separate at $u_v$, write the coefficient as its edge
part plus $(1-d_v)u_v$. At a terminal singleton the existing-label hit
has coefficient $1$. These are exactly the three listed contributions.

To compute their sum, let the occupied edge lengths be $\ell_e$.
Summing endpoint differences over the finite tree gives
\[
 \sum_e\ell_e+\sum_v(d_v-1)u_v=r-a.
\]
Indeed each terminal contributes $1$, the initial point contributes
$-a$, and at a fork the endpoint contributions equal
$-(d_v-1)u_v$. All contributions except the $r$ coincidences are
nonpositive. Their signed sum is $a$, and their total absolute sum is
$r+(r-a)$. If there is no occupied child, the one-slot coefficient is
simply $a$. Insert further labels sequentially and apply the variation
bound to the enlarged occupied tree. The row-sum assertion follows
before taking absolute values from the total signed mass calculation.
Empty covariance steps merely subdivide edges and change none of these
identities.
\end{proof}

\begin{remark}[What a small mass does and does not imply]
The first entry into an empty block of mass $u$ costs a factor $u$.
Likewise departures on an initial interval of length $u$ have total
variation at most a constant, depending on the old leaf count, times
$u$. Subsequent fixed numbers of allocations retain this factor by
Lemma~\ref{lem:found-attachment}. However, the coefficient
$a-kb$ need not be $O(a)$ if $b\gg a$. A small enclosing mass alone
does not supply a factor to every ungrouped term. For two slots in an
empty block, the two different-child coefficient is $a(a-b)$ and the
same-child contribution on a constant test is $ab$. Their signed total
is $a^2$, whereas their absolute sum need only be $O(a)$. All small-mass
savings below use an actual first entry or an exact row collapse.
\end{remark}

\begin{lemma}[Densities of split levels]\label{lem:found-split-density}
Fix a finite observed tree with at most $q$ leaves and allocate at most
$k$ new labels. Apart from coincidences and atoms at specified old forks
or cuts, the split mass of any specified pair involving a new label has
a Lebesgue density with total variation bounded by a constant $C_{q,k}$.
The same statement holds when some old fork masses are themselves
integrated against a measure of bounded density, with the resulting
bound multiplied by that density bound. These assertions concern the
signed attachment measure and its variation, not the size of a
spin observable at a pointwise split level.
\end{lemma}

\begin{proof}
Allocate the specified pair before the other new labels, using equality
of the underlying finite product derivatives to reorder the allocations.
If one label is old, its non-atomic separation from the new label occurs
on an occupied edge and has density of absolute value $1$; there are at
most a number of edges depending on $q$. Old forks give precisely the
listed atoms. If both labels are new, first place one of them. A fork
created by a placement on an edge already has a bounded mass density.
The second label either separates on an edge, attaches at an old fork,
or attaches at this newly created fork. The last coefficient is bounded
by the number of occupied children, because its mass is at most $1$.
It therefore preserves the density bound in that variable. Coincidences
are kept separate. All remaining allocations have bounded total
variation by Lemma~\ref{lem:found-attachment}. Integration over an old
fork variable with bounded density is the same argument followed by
Fubini. The reordering was an exact identity before any absolute values
were taken, so it introduces no cancellation assumption in the final
variation estimate.
\end{proof}

\begin{lemma}[Uniform Taylor bounds]\label{lem:found-taylor}
Under the hypotheses of Proposition~\ref{prop:found-replica}, suppose
$P$ uses at most $q$ labels and $|V(u;\sigma,\tau)|\leq M$. For each
fixed integer $j\geq1$ there is $C_{q,j}<\infty$, independent of the
number of covariance steps, such that
\[
 \left|\frac{d^j}{d\theta^j}\nu_\theta[P]\right|
       \leq C_{q,j}M^j\|P\|_\infty.
\]
More precisely the derivative is an integral over a signed measure of
extended trees with uniformly bounded total variation, of expectations
of $P$ times $j$ covariance insertions, with the factor $2^{-j}$.
H\"older's inequality can therefore be applied to these integrands
under their genuine probability laws. If $P\geq0$, then
\[
 \left|\frac{d}{d\theta}\nu_\theta[P]\right|
        \leq C_{q,1}M\nu_\theta[P],
\]
and hence
\[
 e^{-C_{q,1}M|\theta-\theta'|}\nu_{\theta'}[P]
 \leq\nu_\theta[P]
 \leq e^{C_{q,1}M|\theta-\theta'|}\nu_{\theta'}[P].
\]
The assertions require that the test, including any scalar centers in
it, be fixed during the interpolation. For parameter-dependent tests
the ordinary product-rule derivatives must also be included.
\end{lemma}

\begin{proof}
Iterate \eqref{eq:2}. Affineness makes every covariance insertion fixed,
so the $j$th derivative has $2j$ additional labels and $j$ insertions.
Lemma~\ref{lem:found-attachment} bounds the variation of the resulting
measure by a constant depending only on $q,j$. This proves the first
two assertions, including Taylor's integral remainder. For a nonnegative
old-label test, each extended law has the same old-label marginal.
Consequently $\nu_\theta[|P V_{cd}|]\leq M\nu_\theta[P]$ for every
attachment, giving the differential inequality. Integrating it, or using
the integral form of Gr\"onwall's inequality when the expectation
vanishes, proves the comparison.
\end{proof}

\begin{corollary}[Source derivative bounds]\label{cor:found-source-bounds}
For fixed $n$, every derivative order of a continuation in bounded
feature directions has a bound independent of the mass partition,
uniform over the external source and the admissible bounded covariance
clocks. The bound depends only on the order and the suprema of the
feature directions.
\end{corollary}

\begin{proof}
The first derivative is the conditional feature mean. Differentiate it
with the one-slot version of Proposition~\ref{prop:found-replica}, and
repeat. Every additional derivative introduces one bounded feature and
one summation label. Lemma~\ref{lem:found-attachment} bounds each
successive attachment measure. A continuation beginning at a prescribed
prefix is treated as a rooted hierarchy by adjoining empty earlier
steps; its prefix is a deterministic external source. This makes the
same argument applicable without any lower bound on its first mass.
\end{proof}

\subsection{Cut truncations and conditional couplings}

We record the precise finite-tree facts needed when a free row is
replaced by a row truncated at a cut. A detailed attachment of a new
label $r$ is truncated at $t$ as follows. If it already departs from the
old tree before $t$, nothing is changed. Otherwise retain its old
ancestor through the prescribed revelation at $t$ and then draw its
continuation independently of the old descendants. Thus deeper
placements in the same occupied block at $t$ give the same truncated
law. The old-label marginal is always unchanged.

\begin{lemma}[Variation of a cut layer]\label{lem:found-cut-truncation}
Fix an old tree with at most $q$ leaves in a block of mass $U$, and a
free row label $r$ in that block. Let $U\leq v\leq w\leq\min(2v,1)$.
The difference of its signed row tests truncated at $w$ and at $v$ can,
after exact grouping of the deeper placements, be written as a signed
combination of pairs of planted laws with total coefficient variation
at most $C_qv$. In each pair the old-label tree is the same, and only
the coupling of $r$ to one occupied descendant subtree after the
revelation at $v$ differs. One may take $C_q=12q$.

The statement remains true when the row is multiplied by a bounded
coefficient depending on the fixed old topology, multiplied into the
bound by its supremum. For the exact collapse of deeper placements this
coefficient, and all scalar centers in the row test, must be independent
of $r$'s detailed attachment beyond the truncation cut.
Allocation-order symmetry carries the full topology coefficient with it;
it does not establish the independence required for a free-row collapse.
At each recursion generation, freeze the previously generated topology,
its coefficients, and earlier scalar centers before introducing fresh
slots. A current coefficient depending on the row's deeper placement
must remain explicit until eliminated or estimated. Only after this
row calculation may the enlarged topology be frozen for the next
generation.
\end{lemma}

\begin{proof}
Contributions departing before $v$ are identical and cancel. There are
at most $q$ occupied subblocks at either cut. Inside an occupied subblock
at $w$, every deeper placement gives the same $w$-truncated expectation,
so its signed sum is exactly $w$ by
Lemma~\ref{lem:found-attachment}. These collapsed tails have total
absolute coefficient at most $qw\leq2qv$. Departures on edges between
$v$ and $w$ have total length at most $q(w-v)\leq qv$. Fork atoms in
this layer have total absolute weight at most
$w\sum(d_v-1)\leq2qv$. Treat forks lying exactly on a cut with the
chosen revelation side; the same estimate covers them. Each of these
grouped contributions is the difference between retaining its indicated
coupling after $v$ and making $r$ independent at $v$; the old marginal
is preserved in both laws. Counting the two terms of these pairs gives
at most $10qv$, and $12qv$ allows terminal coincidences with either cut
convention. A terminal coincidence can occur in the layer only if
$w=1$, in which case $v\geq1/2$ and its unit weight is covered by the
same bound. Before that layer it is already included in a collapsed
tail. The asserted restriction on coefficients is exactly what is
needed for the identical-expectation grouping, and no further
independence from the old tree is used.
\end{proof}

\begin{lemma}[Centered coupling difference]\label{lem:found-coupling}
Consider either pair of laws in Lemma~\ref{lem:found-cut-truncation}.
Let $\mathcal G$ be the sigma-field generated by their complete common
prefix through the prescribed revelation at $v$, the realizations of
all old branches outside the affected descendant subtree after $v$,
and the deterministic old topology. No information from the continuation
of $r$ or from the affected old subtree after $v$ is included in
$\mathcal G$. Let $F$ be a function of $r$ and
this exterior information, and let $P$ be a function of the old
subtree and the exterior information. The two laws have the same
conditional marginal law of $F$ given $\mathcal G$ and the same
conditional marginal law of $P$ given $\mathcal G$. Write their common
conditional means as $\bar F,\bar P$. Their difference therefore obeys
\[
 (\mathbf E_1-\mathbf E_0)[FP]
   =(\mathbf E_1-\mathbf E_0)[(F-\bar F)(P-\bar P)].
\]
If $p^{-1}+p'^{-1}=1$, its absolute value is at most
\[
 \sum_{i=0}^1\|F-\bar F\|_{L^p(\mathbf P_i)}
                  \|P-\bar P\|_{L^{p'}(\mathbf P_i)}.
\]
Each centered norm is at most the corresponding norm of the difference
from an independent resampling with the same conditional marginal given
$\mathcal G$. Consequently uniform resampling bounds $A_v,B_v$ for
these two factors give a bound $2C_qvA_vB_v$ for the full cut layer.
\end{lemma}

\begin{proof}
Normalize and integrate every branch not retained in the conditioning.
Normalization of the transition kernels preserves the old subtree law.
It also gives the same single-continuation marginal for $r$, whether
$r$ follows an old branch for longer, coincides with an old terminal
label, or is made independent immediately after $v$. Conditional on
$\mathcal G$, the exterior information can change the distribution of
the shared prefix but not the normalized future kernels in the affected
subtree. This proves the two marginal assertions. Expand
$FP=(F-\bar F)(P-\bar P)+\bar FP+\bar PF-\bar F\bar P$;
the last three terms have equal expectations under the two laws.
H\"older's inequality gives the displayed bound after averaging over
$\mathcal G$. If $F'$ is an independent conditional copy, then
$F-\bar F=\mathbf E[F-F'\mid F,\mathcal G]$, so Jensen's inequality
gives the resampling bound; the argument for $P$ is identical.
Finally integrate against the coefficient measure of
Lemma~\ref{lem:found-cut-truncation}. Conditional bounds uniform in
the realized prefix are not needed: bounds in the displayed
unconditional $L^p$ norms suffice.
\end{proof}

\begin{lemma}[Boundary tests]\label{lem:found-boundary-tests}
Fix a cut $U$ and its old descendant blocks. Let $P$ be a polynomial of
fixed degree in overlaps between different blocks, centered by fixed
scalars, and in products of such overlaps with fresh auxiliary leaves
split at specified cuts. Assume that each overlap involving an affected
old subtree has its other endpoint outside that subtree. Bounded
coefficients may depend on the entire fixed old topology. All scalar
centers may be expectations on an earlier tree, provided they are held
fixed during the current row interpolation.

On resampling the affected subtree while retaining its prefix, only
the overlap factors crossing its boundary vary. If each such factor
has a resampling-difference bound in the required finite $L^p$ norm,
and the other factors have their required finite moment bounds, then
$P$ has the bound obtained by the product rule, with a constant
depending only on its degree, number of monomials, and coefficient bound.
In particular this property is preserved by multiplying the current
target test by centered overlaps to fresh auxiliary leaves, and by
introducing a fresh target block joined to an earlier target by one
overlap. Earlier target blocks may have several incident boundary
overlaps.
\end{lemma}

\begin{proof}
For a monomial use the exact identity
\[
 \prod_{i=1}^d A_i-\prod_{i=1}^d A_i'
    =\sum_{j=1}^d(A_j-A_j')
                   \prod_{i<j}A_i\prod_{i>j}A_i'.
\]
Factors supported in the exterior have $A_j=A_j'$. A fixed scalar center
also cancels from $A_j-A_j'$. Apply H\"older's inequality to each
remaining term, at exponents whose reciprocal sum is the reciprocal of
the desired exponent. Sum over the fixed finite collection of monomials.
This proves the norm transfer without any assertion that internal
overlaps are concentrated. The two stated operations introduce edges
between the current target and another block; they create no edge with
both endpoints in a single descendant block. The same argument
therefore remains applicable after each operation. The coefficient and
centering restrictions ensure that no unrecorded difference term is
introduced.
\end{proof}

\subsection{Passage to bounded monotone paths}

\begin{proposition}[Path limits]\label{prop:found-limits}
The preceding construction and estimates extend to bounded nonnegative
nondecreasing paths $b,h$ by step approximation. The limit is independent
of the approximations, provided the terminal clocks converge and the
before/after clocks at every specified cut converge to the specified
values. Planted expectations of bounded spin tests on a fixed finite
split tree with cuts before or after jumps, fixed-order source derivatives,
the signed calculus, and the mass-integrated pressure and moment formulas
have these limits. The susceptibility inequality holds on the continuous
parts, with nonnegative total increments of $S$ at jumps. For intermediate
cuts and the separate clock contributions inside a filled jump, use
approximations retaining the prescribed joint filling; no independence
from the choice of that filling is asserted. All constants in
the finite-tree variation and Taylor bounds are unchanged.
\end{proposition}

\begin{proof}
We give the details with $n$ fixed; no quantitative estimate uniform in
$n$ is inferred from this approximation argument. A bounded monotone
function has only countably many discontinuities. Choose step
approximations converging at its continuity points and retaining each
required side of the finitely many specified cuts. They converge in
$L^1$, and their endpoint clocks can be chosen to converge as well.
Lemma~\ref{lem:found-pressure} makes the self-corrected root values a
Cauchy sequence, uniformly over every common external source.

The same calculation applies to a continuation beginning at a fixed
cut $u$. Its future increment clocks are the cumulative clocks above
$u$ minus their specified values at that cut. Heat differentiation and
summation by parts give the same interior integral, together with an
initial boundary term. The latter is bounded by a constant depending
on $n$ times the difference of the two clocks at $u$, because all
feature means have norm at most their feature bound. Thus continuation
values converge uniformly in the external source whenever the clocks
at the cut converge; the bound consists of their $L^1$ differences
above $u$, their initial clock differences, and their terminal self
corrections. This also proves uniqueness of these limits.

Corollary~\ref{cor:found-source-bounds} bounds all source derivatives
uniformly over the approximations. Uniform convergence of values and
the bound on the next derivative imply convergence of any fixed
derivative by finite differences: approximate a derivative by a fixed
small difference quotient, pass to the limit in the quotient, and then
let its step tend to zero. Iterating this argument proves the claim
for mixed derivatives as well. In particular the limiting gradient and
Hessian are the conditional mean and weighted covariance limits.

To verify convergence of the probability laws, use independent standard
Brownian drivers for the finitely many features on each branch of the
fixed observed tree. Evaluate them at the deterministic $b,h$ clocks.
Continuity of Brownian paths gives convergence at every cut whose clocks
converge, including its prescribed side. Multiplication of normalized
transition densities telescopes on each edge exactly as in the pinned
calculation: the logarithm of the joint density is a finite sum of
endpoint and fork continuation terms and terminal Hamiltonians, minus
integrals of continuations along the edges in the mass coordinate.
The coefficients of endpoint terms are bounded in terms of the number
of observed leaves.

Here is a uniform integrable bound for this density. For fixed $n$ and
bounded clocks, terminal log partition functions are Lipschitz in the
feature source with fixed constant. The same is true of continuations
by the first derivative formula. Their values at source zero are
nonnegative and bounded by the annealed value, which is at most
$n(b(1-)/4+h(1-)/2)$. It follows that the absolute value of the
telescoped log density is at most
\[
 C_{n,q}\left(1+\sum_{\text{branches and features}}
                      \sup_{0\leq t\leq M}|B_t|\right),
\]
where $M$ bounds all clocks and $q$ is the fixed leaf count. The
exponential of this expression is integrable for a finite collection
of Brownian motions. Continuation convergence and dominated convergence
in the mass integrals therefore give almost sure convergence of the
densities and convergence in $L^1$. The same domination, together with
the source derivative bounds, applies to the differentiated observables
under consideration. This proves convergence of positive planted
expectations and their fixed-order source versions.

The attachment measure consists of finitely many fork and coincidence
atoms and Lebesgue measures on finitely many occupied edges, with
uniform total variation. At a Lebesgue placement, the clocks converge
outside a countable exceptional set. At an existing fork they converge
on the specified side by construction. Dominated convergence and
Lemma~\ref{lem:found-attachment} therefore pass every fixed finite
sequence of signed allocations to the limit. The same argument on an
affine interpolation segment passes fixed-order differentiation and
Taylor's integral remainder, since their variation bounds are uniform.
Exact permutation and row-sum identities survive because they held
before passage to the limit.

Finally $|L_u|\leq3/2$ and $0\leq W_u\leq4$. Thus the moment estimate
passes first for bounded deterministic weights and then for all
$L^2$ weights. For the clock identity, choose refinements that retain the prescribed
joint filling of every jump whose total clock size exceeds a cutoff,
and then let the cutoff decrease to zero. After the finitely many retained
jumps are separated, remove the remaining jumps temporarily. Their total
clock size tends to zero; the continuation and source-derivative estimates
above control this perturbation. The resulting clocks are continuous on
each interval between retained jumps, so their monotonicity gives uniform
clock convergence there. These estimates and Brownian continuity give
convergence of the quadratic-variation integrands, including within the
vanishing increments of the refinements, and hence of their Stieltjes
integrals. The uniform integrand bounds control the direct contribution
of the omitted jumps as well. At each retained jump the
prescribed joint filling fixes the separate clock integrals. Their sum
is the total nonnegative increment of $S$, independently of the filling.
Discarding this increment gives the susceptibility inequality on the
continuous portions. This proves the stated clock assertions without
claiming that the separate jump contributions are independent of their
joint filling.
\end{proof}

Henceforth all path quantities refer to these limits. The signed sums
are integrals of ordinary finite-tree expectations against the
attachment measures above. In particular all subsequent fixed-order
Taylor estimates have constants independent of how finely a path is
approximated.
